\documentclass[aps,preprint]{revtex4}%
\usepackage{amsfonts}
\usepackage{amsmath}
\usepackage{amssymb}
\usepackage{graphicx}%
\setcounter{MaxMatrixCols}{30}
%TCIDATA{OutputFilter=latex2.dll}
%TCIDATA{Version=5.50.0.2953}
%TCIDATA{CSTFile=revtex4.cst}
%TCIDATA{Created=Friday, July 18, 2008 09:51:28}
%TCIDATA{LastRevised=Friday, August 08, 2008 14:06:56}
%TCIDATA{<META NAME="GraphicsSave" CONTENT="32">}
%TCIDATA{<META NAME="SaveForMode" CONTENT="1">}
%TCIDATA{BibliographyScheme=Manual}
%TCIDATA{<META NAME="DocumentShell" CONTENT="Articles\SW\REVTeX 4">}
%BeginMSIPreambleData
\providecommand{\U}[1]{\protect\rule{.1in}{.1in}}
%EndMSIPreambleData
\newtheorem{theorem}{Theorem}
\newtheorem{acknowledgement}[theorem]{Acknowledgement}
\newtheorem{algorithm}[theorem]{Algorithm}
\newtheorem{axiom}[theorem]{Axiom}
\newtheorem{claim}[theorem]{Claim}
\newtheorem{conclusion}[theorem]{Conclusion}
\newtheorem{condition}[theorem]{Condition}
\newtheorem{conjecture}[theorem]{Conjecture}
\newtheorem{corollary}[theorem]{Corollary}
\newtheorem{criterion}[theorem]{Criterion}
\newtheorem{definition}[theorem]{Definition}
\newtheorem{example}[theorem]{Example}
\newtheorem{exercise}[theorem]{Exercise}
\newtheorem{lemma}[theorem]{Lemma}
\newtheorem{notation}[theorem]{Notation}
\newtheorem{problem}[theorem]{Problem}
\newtheorem{proposition}[theorem]{Proposition}
\newtheorem{remark}[theorem]{Remark}
\newtheorem{solution}[theorem]{Solution}
\newtheorem{summary}[theorem]{Summary}
\newenvironment{proof}[1][Proof]{\noindent\textbf{#1.} }{\ \rule{0.5em}{0.5em}}
\begin{document}
\preprint{HEP/123-qed}
\title[Short title for running header]{Long title that appears on the title page}
\author{First Author}
\affiliation{My Institution}
\author{Second Author}
\affiliation{My Institution}
\author{Third Author}
\affiliation{Other Institution}
\keywords{one two three}
\pacs{PACS number}

\begin{abstract}
Shell document for REV\TeX{} 4.

\end{abstract}
\volumeyear{year}
\volumenumber{number}
\issuenumber{number}
\eid{identifier}
\date[Date text]{date}
\received[Received text]{date}

\revised[Revised text]{date}

\accepted[Accepted text]{date}

\published[Published text]{date}

\startpage{101}
\endpage{102}
\maketitle
\tableofcontents


The position-spin field operators are defined by
\begin{equation}
\Phi\left(  \mathbf{x}\right)  ^{\dagger}=\sum_{1\leq j\leq r_{1}}\varphi
_{j}\left(  \mathbf{x}\right)  a_{j}^{\dagger}%
\end{equation}
giving the inverse relationship%
\begin{equation}
\int\varphi_{k}^{\ast}\left(  \mathbf{x}\right)  \Phi\left(  \mathbf{x}%
\right)  ^{\dagger}d\mathbf{x}=\sum_{1\leq j\leq r_{1}}\int\varphi_{k}^{\ast
}\left(  \mathbf{x}\right)  \varphi_{j}\left(  \mathbf{x}\right)
d\mathbf{x}a_{j}^{\dagger}=a_{k}^{\dagger}.
\end{equation}
Thus any local operator can be expressed as%
\begin{equation}
\int\lambda\left(  \mathbf{x}\right)  \Phi\left(  \mathbf{x}\right)
^{\dagger}\Phi\left(  \mathbf{x}\right)  d\mathbf{x}=\sum_{1\leq j,k\leq
r_{1}}\left\{  \int\lambda\left(  \mathbf{x}\right)  \varphi_{j}\left(
\mathbf{x}\right)  \varphi_{k}^{\ast}\left(  \mathbf{x}\right)  d\mathbf{x}%
\right\}  a_{j}^{\dagger}a_{k}=\sum_{1\leq j,k\leq r_{1}}\lambda_{kj}%
a_{j}^{\dagger}a_{k}%
\end{equation}
where there are certain conditions on $\left\{  \lambda_{kj}\right\}  $ apart
from $\lambda_{kj}=\lambda_{jk}^{\ast}.$%
\begin{align}
\lambda & =\sum_{1\leq j,k\leq r_{1}}\left\vert \varphi_{j}\right\rangle
\lambda_{jk}\left\langle \varphi_{k}\right\vert =\int\left\vert \mathbf{x}%
\right\rangle \lambda\left(  \mathbf{x,x}^{\prime}\right)  \left\langle
\mathbf{x}^{\prime}\right\vert d\mathbf{x}d\mathbf{x}^{\prime}=\int\left\vert
\mathbf{x}\right\rangle \lambda\left(  \mathbf{x}\right)  \delta\left(
\mathbf{x-x}^{\prime}\right)  \left\langle \mathbf{x}^{\prime}\right\vert
d\mathbf{x}d\mathbf{x}^{\prime}\nonumber\\
& \Rightarrow\left\langle \mathbf{x}\left\vert \lambda\right.  \mathbf{x}%
^{\prime}\right\rangle =\left\langle \mathbf{x}\left\vert \sum_{1\leq j,k\leq
r_{1}}\left\vert \varphi_{j}\right\rangle \lambda_{jk}\left\langle \varphi
_{k}\right\vert \right.  \mathbf{x}^{\prime}\right\rangle =\sum_{1\leq j,k\leq
r_{1}}\left\langle \mathbf{x}\left\vert \varphi_{j}\right.  \right\rangle
\lambda_{jk}\left\langle \varphi_{k}\left\vert \mathbf{x}^{\prime}\right.
\right\rangle \nonumber\\
& =\sum_{1\leq j,k\leq r_{1}}\left\langle \mathbf{x}\left\vert \varphi
_{j}\right.  \right\rangle \lambda_{jk}\left\langle \varphi_{k}\left\vert
\mathbf{x}^{\prime}\right.  \right\rangle =\lambda\left(  \mathbf{x}\right)
\delta\left(  \mathbf{x-x}^{\prime}\right)
\end{align}


$\lceil$%
\[
\int\lambda\left(  \mathbf{x}\right)  \varphi_{j}\left(  \mathbf{x}\right)
\varphi_{k}^{\ast}\left(  \mathbf{x}\right)  d\mathbf{x=}\int\lambda\left(
\mathbf{x,x}^{\prime}\right)  \varphi_{j}\left(  \mathbf{x}^{\prime}\right)
\varphi_{k}^{\ast}\left(  \mathbf{x}\right)  d\mathbf{x}d\mathbf{x}^{\prime
}=\int\eta\left(  \mathbf{x}\right)  \varphi_{k}^{\ast}\left(  \mathbf{x}%
\right)  d\mathbf{x}%
\]
where%
\[
\eta_{j}\left(  \mathbf{x}\right)  =\int\lambda\left(  \mathbf{x,x}^{\prime
}\right)  \varphi_{j}\left(  \mathbf{x}^{\prime}\right)  d\mathbf{x}^{\prime}.
\]
Let
\[
\zeta_{j}\left(  \mathbf{x}\right)  =\left\{
\begin{array}
[c]{c}%
\frac{\eta\left(  \mathbf{x}\right)  }{\varphi_{j}\left(  \mathbf{x}\right)
};\varphi_{j}\left(  \mathbf{x}\right)  \neq0\\
0;\varphi_{j}\left(  \mathbf{x}\right)  =0
\end{array}
\right.
\]
then
\[
\int\lambda\left(  \mathbf{x,x}^{\prime}\right)  \varphi_{j}\left(
\mathbf{x}^{\prime}\right)  \varphi_{k}^{\ast}\left(  \mathbf{x}\right)
d\mathbf{x}d\mathbf{x}^{\prime}=\int\zeta_{j}\left(  \mathbf{x}\right)
\varphi_{j}\left(  \mathbf{x}\right)  \varphi_{k}^{\ast}\left(  \mathbf{x}%
\right)  d\mathbf{x.\rfloor}%
\]
The effect of the components of the density operator is given by
\begin{align*}
\Phi\left(  \mathbf{x}\right)  ^{\dagger}\Phi\left(  \mathbf{x}\right)
\left\vert \psi_{1}\cdots\psi_{N}\right\rangle  &  =\sum_{1\leq j,k\leq r_{1}%
}\varphi_{j}\left(  \mathbf{x}\right)  \varphi_{k}^{\ast}\left(
\mathbf{x}\right)  a_{j}^{\dagger}a_{k}\left\vert \psi_{1}\cdots\psi
_{N}\right\rangle \\
&  =\sum_{1\leq h,h^{\prime}\leq N}\varphi_{h}\left(  \mathbf{x}\right)
\varphi_{h^{\prime}}^{\ast}\left(  \mathbf{x}\right)  a_{h}^{\dagger
}a_{h^{\prime}}\left\vert \psi_{1}\cdots\psi_{N}\right\rangle +\sum
_{\substack{1\leq h\leq N\\N+1\leq p\leq r_{1}}}\varphi_{p}\left(
\mathbf{x}\right)  \varphi_{h}^{\ast}\left(  \mathbf{x}\right)  a_{p}%
^{\dagger}a_{h}\left\vert \psi_{1}\cdots\psi_{N}\right\rangle \\
&  =\left\{  \sum_{1\leq h\leq N}\varphi_{h}\left(  \mathbf{x}\right)
\varphi_{h}^{\ast}\left(  \mathbf{x}\right)  \right\}  \left\vert \psi
_{1}\cdots\psi_{N}\right\rangle +\sum_{\substack{1\leq h\leq N\\N+1\leq p\leq
r_{1}}}\varphi_{p}\left(  \mathbf{x}\right)  \varphi_{h}^{\ast}\left(
\mathbf{x}\right)  \left\vert \psi_{1}\cdots\psi_{p}\widehat{\psi}_{h}%
\cdots\psi_{N}\right\rangle
\end{align*}
leading to%
\begin{align*}
\int\lambda\left(  \mathbf{x}\right)  \Phi\left(  \mathbf{x}\right)
^{\dagger}\Phi\left(  \mathbf{x}\right)  d\mathbf{x}\left\vert \psi_{1}%
\cdots\psi_{N}\right\rangle  &  =\sum_{1\leq j,k\leq r_{1}}\lambda_{kj}%
a_{j}^{\dagger}a_{k}\left\vert \psi_{1}\cdots\psi_{N}\right\rangle \\
&  =\left\{  \sum_{1\leq h\leq N}\lambda_{hh}\right\}  \left\vert \psi
_{1}\cdots\psi_{N}\right\rangle +\sum_{\substack{1\leq h\leq N\\N+1\leq p\leq
r_{1}}}\lambda_{ph}\left\vert \psi_{1}\cdots\psi_{p}\widehat{\psi}_{h}%
\cdots\psi_{N}\right\rangle .
\end{align*}
Thus%
\[
\int\lambda\left(  \mathbf{x}\right)  \Phi\left(  \mathbf{x}\right)
^{\dagger}\Phi\left(  \mathbf{x}\right)  d\mathbf{x}\left\vert \psi_{1}%
\cdots\psi_{N}\right\rangle =0\Leftrightarrow\sum_{1\leq h\leq N}\lambda
_{hh}=0\text{ and }\lambda_{ph},\lambda_{hp}=0;1\leq h\leq N,N+1\leq p\leq
r_{1}%
\]
and the local potential operator can be expanded as
\begin{align*}
\Omega\left(  \lambda\right)   &  =\sum_{1\leq h\neq h^{\prime}\leq N}%
\lambda_{hh^{\prime}}a_{h}^{\dagger}a_{h^{\prime}}+\sum_{N+1\leq p,p^{\prime
}\leq N}\lambda_{pp^{\prime}}a_{p}^{\dagger}a_{p^{\prime}}\\
&  =\int\left\{  \sum_{1\leq h\neq h^{\prime}\leq N}\varphi_{h}^{\ast}\left(
\mathbf{x}\right)  \lambda_{hh^{\prime}}\varphi_{h^{\prime}}\left(
\mathbf{x}^{\prime}\right)  +\sum_{N+1\leq p,p^{\prime}\leq N}\varphi
_{p}^{\ast}\left(  \mathbf{x}\right)  \lambda_{pp^{\prime}}\varphi_{p^{\prime
}}\left(  \mathbf{x}^{\prime}\right)  \right\}  \Phi\left(  \mathbf{x}\right)
^{\dagger}\Phi\left(  \mathbf{x}^{\prime}\right)  d\mathbf{x}d\mathbf{x}%
^{\prime},
\end{align*}
showing that%
\[
\lambda\left(  \mathbf{x}\right)  =\lambda\left(  \mathbf{x}\right)
\delta\left(  \mathbf{x-x}^{\prime}\right)  =\sum_{1\leq h\neq h^{\prime}\leq
N}\varphi_{h}^{\ast}\left(  \mathbf{x}\right)  \lambda_{hh^{\prime}}%
\varphi_{h^{\prime}}\left(  \mathbf{x}^{\prime}\right)  +\sum_{N+1\leq
p,p^{\prime}\leq N}\varphi_{p}^{\ast}\left(  \mathbf{x}\right)  \lambda
_{pp^{\prime}}\varphi_{p^{\prime}}\left(  \mathbf{x}^{\prime}\right)  .
\]


\begin{theorem}
Let
\begin{equation}
\Phi\left(  \mathbf{x}\right)  ^{\dagger}\Phi\left(  \mathbf{x}\right)
\left\vert \psi_{1}\cdots\psi_{N}\right\rangle \overset{a.e}{\neq
}0;\,\mathbf{x\in\mathbf{\Delta}\left(  \psi_{1}\cdots\psi_{N}\right)
\subseteq}\mathbb{X}%
\end{equation}
then
\begin{equation}
\Omega\left(  \lambda\right)  \left\vert \psi_{1}\cdots\psi_{N}\right\rangle
=\left\vert \Omega\left(  \lambda\right)  \psi_{1}\cdots\psi_{N}\right\rangle
+\cdots+\left\vert \psi_{1}\cdots\Omega\left(  \lambda\right)  \psi
_{N}\right\rangle =0\Leftrightarrow\lambda\overset{a.e}{=}0.
\end{equation}
(By a.e. we mean almost everywhere in $\mathbf{\mathbf{\Delta}\left(  \psi
_{1}\cdots\psi_{N}\right)  )}$
\end{theorem}

\begin{proof}%
\begin{align}
&  \Phi\left(  \mathbf{x}\right)  ^{\dagger}\Phi\left(  \mathbf{x}\right)
\left\vert \psi_{1}\cdots\psi_{N}\right\rangle \overset{a.e}{\neq
}0;\,\mathbf{x\in}\Delta\left(  \psi_{1}\cdots\psi_{N}\right)
\mathbf{\subseteq}\mathbb{X}\nonumber\\
&  \left.  \Rightarrow\right.  \left\langle \left.  \mathbf{x}_{1}%
\cdots\mathbf{x}_{N}\right\vert \Phi\left(  \mathbf{x}\right)  ^{\dagger}%
\Phi\left(  \mathbf{x}\right)  \psi_{1}\cdots\psi_{N}\right\rangle
\overset{a.e}{\neq}0
\end{align}
If
\begin{equation}
\Phi\left(  \mathbf{x}\right)  ^{\dagger}\Phi\left(  \mathbf{x}\right)
\left\vert \psi_{1}\cdots\psi_{N}\right\rangle \neq0
\end{equation}
then
\begin{equation}
\left\langle \left.  \mathbf{xx}_{2}\cdots\mathbf{x}_{N}\right\vert
\Phi\left(  \mathbf{x}\right)  ^{\dagger}\Phi\left(  \mathbf{x}\right)
\psi_{1}\cdots\psi_{N}\right\rangle =\left\langle \left.  \mathbf{xx}%
_{2}\cdots\mathbf{x}_{N}\right\vert \psi_{1}\cdots\psi_{N}\right\rangle
\overset{a.e}{\neq}0.
\end{equation}
Hence in particular
\begin{equation}
\left\langle \left.  \mathbf{x}_{1}\cdots\mathbf{x}_{N}\right\vert \psi
_{1}\cdots\psi_{N}\right\rangle \overset{a.e}{\neq}0;\quad\mathbf{x}_{j}%
\in\Delta\left(  \psi_{1}\cdots\psi_{N}\right)  \mathbf{;}1\mathbf{\leq}j\leq
N.
\end{equation}
%

\begin{equation}
\gamma\left(  \mathbf{x}\right)  =\left\langle \psi_{1}\cdots\psi_{N}\left.
\Phi\left(  \mathbf{x}\right)  ^{\dagger}\Phi\left(  \mathbf{x}\right)
\right\vert \psi_{1}\cdots\psi_{N}\right\rangle \overset{a.e}{\neq
}0\Longleftrightarrow\Phi\left(  \mathbf{x}\right)  ^{\dagger}\Phi\left(
\mathbf{x}\right)  \left\vert \psi_{1}\cdots\psi_{N}\right\rangle
\overset{a.e}{\neq}0
\end{equation}
In the case $N=1$ one has%
\begin{align}
\int\lambda\left(  \mathbf{x}\right)  \Phi\left(  \mathbf{x}\right)
^{\dagger}\Phi\left(  \mathbf{x}\right)  \left\vert \psi_{j}\right\rangle
d\mathbf{x}  &  =\int\lambda\left(  \mathbf{x}\right)  \Phi\left(
\mathbf{x}\right)  ^{\dagger}\Phi\left(  \mathbf{x}\right)  \left\vert
\mathbf{x}^{\prime}\right\rangle \left\langle \left.  \mathbf{x}^{\prime
}\right\vert \psi_{j}\right\rangle d\mathbf{x}d\mathbf{x}^{\prime}=\int
\lambda\left(  \mathbf{x}\right)  \left\vert \mathbf{x}\right\rangle
\left\langle \left.  \mathbf{x}\right\vert \psi_{j}\right\rangle
d\mathbf{x}=0\nonumber\\
&  \equiv\int\lambda\left(  \mathbf{x}\right)  \psi_{j}\left(  \mathbf{x}%
\right)  \left\vert \mathbf{x}\right\rangle d\mathbf{x=0\Rightarrow}%
\lambda\left(  \mathbf{x}\right)  \psi_{j}\left(  \mathbf{x}\right)
\overset{a.e}{=}0\Rightarrow\lambda\left(  \mathbf{x}\right)  \overset{a.e}%
{=}0.
\end{align}
In the case $N=2$
\begin{equation}
\left\langle \left.  \mathbf{xx}^{\prime}\right\vert \psi_{1}\psi
_{2}\right\rangle =\frac{1}{\sqrt{2}}\left\{  \psi_{1}\left(  \mathbf{x}%
\right)  \psi_{2}\left(  \mathbf{x}^{\prime}\right)  -\psi_{2}\left(
\mathbf{x}\right)  \psi_{1}\left(  \mathbf{x}^{\prime}\right)  \right\}
\end{equation}
and%
\begin{align}
&  \left\langle \left.  \mathbf{xx}^{\prime}\right\vert \Omega\left(
\lambda\right)  \psi_{1}\psi_{2}\right\rangle +\left\langle \left.
\mathbf{xx}^{\prime}\right\vert \psi_{1}\Omega\left(  \lambda\right)  \psi
_{2}\right\rangle \nonumber\\
&  =\frac{1}{\sqrt{2}}\left\{  \lambda\left(  \mathbf{x}\right)  \psi
_{1}\left(  \mathbf{x}\right)  \psi_{2}\left(  \mathbf{x}^{\prime}\right)
-\psi_{2}\left(  \mathbf{x}\right)  \lambda\left(  \mathbf{x}^{\prime}\right)
\psi_{1}\left(  \mathbf{x}^{\prime}\right)  +\psi_{1}\left(  \mathbf{x}%
\right)  \lambda\left(  \mathbf{x}^{\prime}\right)  \psi_{2}\left(
\mathbf{x}^{\prime}\right)  -\lambda\left(  \mathbf{x}\right)  \psi_{2}\left(
\mathbf{x}\right)  \psi_{1}\left(  \mathbf{x}^{\prime}\right)  \right\}
\nonumber\\
&  =\frac{1}{\sqrt{2}}\left\{  \lambda\left(  \mathbf{x}\right)  \left\{
\psi_{1}\left(  \mathbf{x}\right)  \psi_{2}\left(  \mathbf{x}^{\prime}\right)
-\psi_{2}\left(  \mathbf{x}\right)  \psi_{1}\left(  \mathbf{x}^{\prime
}\right)  \right\}  +\lambda\left(  \mathbf{x}^{\prime}\right)  \left\{
\psi_{1}\left(  \mathbf{x}\right)  \psi_{2}\left(  \mathbf{x}^{\prime}\right)
-\psi_{2}\left(  \mathbf{x}\right)  \psi_{1}\left(  \mathbf{x}^{\prime
}\right)  \right\}  \right\} \nonumber\\
&  =\frac{1}{\sqrt{2}}\left\{  \lambda\left(  \mathbf{x}\right)
+\lambda\left(  \mathbf{x}^{\prime}\right)  \right\}  \left\{  \psi_{1}\left(
\mathbf{x}\right)  \psi_{2}\left(  \mathbf{x}^{\prime}\right)  -\psi
_{2}\left(  \mathbf{x}\right)  \psi_{1}\left(  \mathbf{x}^{\prime}\right)
\right\}  =\left\{  \lambda\left(  \mathbf{x}\right)  +\lambda\left(
\mathbf{x}^{\prime}\right)  \right\}  \left\langle \left.  \mathbf{xx}%
^{\prime}\right\vert \psi_{1}\psi_{2}\right\rangle .
\end{align}
Hence if
\begin{equation}
\left\langle \left.  \mathbf{xx}^{\prime}\right\vert \Omega\left(
\lambda\right)  \psi_{1}\psi_{2}\right\rangle +\left\langle \left.
\mathbf{xx}^{\prime}\right\vert \psi_{1}\Omega\left(  \lambda\right)  \psi
_{2}\right\rangle \overset{ae}{=}0
\end{equation}
then either
\begin{equation}
\left\langle \left.  \mathbf{xx}^{\prime}\right\vert \psi_{1}\psi
_{2}\right\rangle \overset{ae}{=}0
\end{equation}
or
\begin{equation}
\lambda\left(  \mathbf{x}\right)  +\lambda\left(  \mathbf{x}^{\prime}\right)
\overset{ae}{=}0\quad\mathbf{x\neq x}^{\prime}.
\end{equation}
If
\begin{equation}
\left\langle \left.  \mathbf{xx}^{\prime}\right\vert \psi_{1}\psi
_{2}\right\rangle \overset{ae}{=}0
\end{equation}
then
\begin{equation}
\gamma\left(  \mathbf{x}\right)  =2\int\left\langle \psi_{1}\psi_{2}\left\vert
\mathbf{xx}^{\prime}\right.  \right\rangle \left\langle \left.  \mathbf{xx}%
^{\prime}\right\vert \psi_{1}\psi_{2}\right\rangle d\mathbf{x}^{\prime
}\overset{ae}{=}0
\end{equation}
which is a contradiction. Thus $\left\langle \left.  \mathbf{uv}\right\vert
\psi_{1}\psi_{2}\right\rangle \overset{ae}{\neq}0$ for any $\mathbf{u\neq
v}\in\Delta\left(  \psi_{1}\psi_{2}\right)  $ and implies that
\begin{equation}
\lambda\left(  \mathbf{u}\right)  +\lambda\left(  \mathbf{v}\right)
\overset{a.e}{=}0\quad\mathbf{u\neq v}\in\Delta\left(  \psi_{1}\psi
_{2}\right)
\end{equation}
choosing pairs $\left\{  \mathbf{x,x}^{\prime}\right\}  ,\left\{
\mathbf{x}^{\prime}\mathbf{,x}^{\prime\prime}\right\}  $ and $\left\{
\mathbf{x,x}^{\prime\prime}\right\}  $
\begin{align}
&  \lambda\left(  \mathbf{x}\right)  =-\lambda\left(  \mathbf{x}^{\prime
}\right)  =-\lambda\left(  \mathbf{x}^{\prime\prime}\right) \nonumber\\
&  \lambda\left(  \mathbf{x}^{\prime}\right)  =-\lambda\left(  \mathbf{x}%
^{\prime\prime}\right)  \Rightarrow\lambda\left(  \mathbf{x}^{\prime}\right)
=-\lambda\left(  \mathbf{x}^{\prime}\right)
\end{align}
and hence%
\begin{equation}
\lambda\left(  \mathbf{x}^{\prime}\right)  \overset{ae}{=}0.
\end{equation}
In the case $N=3$%
\begin{equation}
\lambda\left(  \mathbf{x}_{1}\right)  +\lambda\left(  \mathbf{x}_{2}\right)
+\lambda\left(  \mathbf{x}_{3}\right)  \overset{ae}{=}0;\mathbf{x}%
_{1}\mathbf{\neq\mathbf{x}_{2}\neq x}_{3}%
\end{equation}
and we can choose triples $\left\{  \mathbf{x,x}^{\prime},\mathbf{x}%
^{\prime\prime}\right\}  ,\left\{  \mathbf{x,x}^{\prime}\mathbf{,x}%
^{\prime\prime\prime}\right\}  ,\left\{  \mathbf{x,\mathbf{x}^{\prime\prime
},x}^{\prime\prime\prime}\right\}  $and$\left\{  \mathbf{x}^{\prime
}\mathbf{,x}^{\prime\prime},\mathbf{x}^{\prime\prime\prime}\right\}  $ so that%
\begin{equation}
\mathbb{A}\mathbf{\lambda=}\left(
\begin{array}
[c]{cccc}%
0 & 1 & 1 & 1\\
1 & 0 & 1 & 1\\
1 & 1 & 0 & 1\\
1 & 1 & 1 & 0
\end{array}
\right)  \left(
\begin{array}
[c]{c}%
\lambda\left(  \mathbf{x}\right) \\
\lambda\left(  \mathbf{x}^{\prime}\right) \\
\lambda\left(  \mathbf{x}^{\prime\prime}\right) \\
\lambda\left(  \mathbf{x}^{\prime\prime\prime}\right)
\end{array}
\right)  =\left(
\begin{array}
[c]{c}%
0\\
0\\
0\\
0
\end{array}
\right)
\end{equation}
and $\det$ $\mathbb{A}$ $\neq0$ Hence $\mathbf{\lambda=0}$ In the general
case
\begin{equation}
\lambda\left(  \mathbf{x}_{1}\right)  +\cdots+\lambda\left(  \mathbf{x}%
_{N}\right)  \overset{ae}{=}0;\mathbf{x}_{1}\mathbf{\neq\cdots\neq x}_{N}%
\end{equation}
and we can choose $N+1$ points $\left\{  \mathbf{y}_{1},\cdots,\mathbf{y}%
_{N+1}\right\}  $ to obtain
\begin{equation}
\mathbb{A}\mathbf{\lambda=}\left(
\begin{array}
[c]{cccc}%
0 & 1 & \cdots & 1\\
1 & 0 & \cdots & 1\\
\vdots & \vdots & 0 & \vdots\\
1 & \cdots & 1 & 0
\end{array}
\right)  \left(
\begin{array}
[c]{c}%
\lambda\left(  \mathbf{y}_{1}\right) \\
\lambda\left(  \mathbf{y}_{2}\right) \\
\vdots\\
\lambda\left(  \mathbf{y}_{N+1}\right)
\end{array}
\right)  =\left(
\begin{array}
[c]{c}%
0\\
0\\
\vdots\\
0
\end{array}
\right)  .
\end{equation}
One can see that the columns of $\mathbb{A}$ for linearly independent vectors
thus $\mathbb{A}$ is nonsingular.

\begin{corollary}
$\left\{  \Phi\left(  \mathbf{x}\right)  ^{\dagger}\Phi\left(  \mathbf{x}%
\right)  \left\vert \psi_{1}\cdots\psi_{N}\right\rangle ;\,\mathbf{x\in
\mathbf{\Delta}\left(  \psi_{1}\cdots\psi_{N}\right)  \subseteq}%
\mathbb{X}\right\}  $ is a generalized basis for $W\subset\mathcal{H}^{N}$
\end{corollary}

\begin{corollary}
$\left\{  \Omega\left(  \lambda\right)  \left\vert \psi_{j}\right\rangle
=0;\,1\leq j\leq N\right\}  \Rightarrow?$
\end{corollary}
\end{proof}


\end{document}